% !TEX root = ../../ctfp-print.tex

\lettrine[lhang=0.17]{在}{数学中，我们有多种方式表达一个事物与另一个事物}「相似」。最严格的相似性是\newterm{等价性（equality）}。当两个事物无法相互区分时，它们就是相等的。在任意可想象的上下文中都可以互相替换。例如，当我们讨论交换图时，是否注意到每次都使用了\newterm{态射的等价性}？这是因为态射构成集合（hom-set），而集合元素可以通过等价性进行比较。

但等价性往往过于严格。在实际应用中存在许多本质上相同但不满足严格等价的例子。例如，有序对类型\code{(Bool, Char)}与\code{(Char, Bool)}虽然不严格相等，但我们知道它们包含完全相同的信息。这种概念最佳体现了两个类型之间的\newterm{同构（isomorphism）}——一种可逆的态射。作为态射，它保持结构；而「同构」意味着通过往返映射可以回到起点，无论从哪一侧开始。对于有序对，这个同构被称为\code{swap}：

\src{snippet01}
\code{swap}恰好是自身的逆映射。

\section{伴随与单位/余单位对}

当我们说两个范畴是同构时，我们通过范畴间的映射（即函子）来表达。我们希望定义：当存在从范畴$\cat{C}$到$\cat{D}$的函子$R$（「右」）且该函子可逆时，称$\cat{C}$与$\cat{D}$同构。换句话说，存在另一个从$\cat{D}$返回$\cat{C}$的函子$L$（「左」），使得与$R$复合后等于恒等函子$I$。这里有两种可能的复合方式：$R \circ L$和$L \circ R$；以及两个恒等函子：分别位于$\cat{C}$和$\cat{D}$中的恒等函子。

\begin{figure}[H]
  \centering
  \includegraphics[width=0.5\textwidth]{images/adj-1.jpg}
\end{figure}

\noindent
但关键问题在于：两个函子\emph{相等}意味着什么？我们如何理解以下等式：
$$R \circ L = I_{\cat{D}}$$
或这个等式：
$$L \circ R = I_{\cat{C}}$$
我们可以尝试用对象的等价性来定义函子等价性：当作用于相等对象时，两个函子应产生相等的对象。但在任意范畴中，我们通常没有对象等价性的概念——这本就不是范畴定义的一部分。（深入探讨「等价性本质」这个问题，最终会导向同伦类型论。）

你可能会争辩说：函子\emph{是}范畴范畴中的态射，因此应该可以比较等价性。确实如此，只要我们讨论的是小范畴（对象构成集合），就可以用集合元素的等价性来比较对象。

然而请注意，$\Cat$本质上是一个$\cat{2}$-范畴。在$\cat{2}$-范畴中，Hom集具有额外结构——存在作用于1-态射之间的2-态射。在$\Cat$中，1-态射是函子，2-态射是自然变换。因此，当我们讨论函子时，更自然（不得不使用这个双关！）地采用自然同构替代等价性。

因此，相较于范畴同构，更合理的广义概念是\newterm{等价性（equivalence）}。若能找到两个双向对应的函子，使得其复合（任一方向）都\newterm{自然同构}于恒等函子，则称范畴$\cat{C}$与$\cat{D}$\emph{等价}。换句话说，存在从复合$R \circ L$到恒等函子$I_{\cat{D}}$的双向自然变换，以及从$L \circ R$到$I_{\cat{C}}$的双向自然变换。

伴随关系比等价性更弱，因为它不要求两个函子的复合\emph{同构}于恒等函子。相反，它规定存在从$I_{\cat{D}}$到$R \circ L$的\newterm{单向}自然变换，以及从$L \circ R$到$I_{\cat{C}}$的另一自然变换。这两个自然变换的签名如下：
\begin{gather*}
  \eta \Colon I_{\cat{D}} \to R \circ L \\
  \varepsilon \Colon L \circ R \to I_{\cat{C}}
\end{gather*}
$\eta$称为伴随的\newterm{单位}，$\varepsilon$称为\newterm{余单位}。

注意这两个定义的不对称性。一般来说，不存在另外两个映射：
\begin{gather*}
  R \circ L \to I_{\cat{D}} \quad\quad\text{不一定存在} \\
  I_{\cat{C}} \to L \circ R \quad\quad\text{不一定存在}
\end{gather*}
由于这种不对称性，函子$L$被称为函子$R$的\newterm{左伴随}，而$R$则是$L$的右伴随。（当然，「左」和「右」的意义取决于你绘制图表的方式。）

伴随关系的紧凑记法为：
$$L \dashv R$$
为了更好地理解伴随关系，让我们更详细地分析单位和余单位。

\begin{figure}[H]
  \centering
  \includegraphics[width=0.5\textwidth]{images/adj-unit.jpg}
\end{figure}

\noindent
从单位开始分析。作为自然变换，它是态射的族。给定$\cat{D}$中的对象$d$，$\eta$的分量是从$I d$（等于$d$）到$(R \circ L) d$的态射；在图中记作$d'$：
$$\eta_d \Colon d \to (R \circ L) d$$
注意复合$R \circ L$是$\cat{D}$上的自函子。

这个等式告诉我们：可以选择$\cat{D}$中任意对象$d$作为起点，通过往返函子$R \circ L$确定目标对象$d'$，然后发射箭头——即态射$\eta_d$——指向目标。

\begin{figure}[H]
  \centering
  \includegraphics[width=0.5\textwidth]{images/adj-counit.jpg}
\end{figure}

\noindent
同理，余单位$\varepsilon$的分量可描述为：
$$\varepsilon_{c} \Colon (L \circ R) c \to c$$
它表明我们可以选择$\cat{C}$中任意对象$c$作为目标，通过往返函子$L \circ R$确定源对象$c' = (L \circ R)c$，然后从源对象发射箭头——即态射$\varepsilon_c$——指向目标。

另一种看待单位和余单位的视角是：单位允许我们在任何可以插入$\cat{D}$恒等函子的地方\emph{引入}复合$R \circ L$；而余单位允许我们\emph{消除}复合$L \circ R$，将其替换为$\cat{C}$的恒等函子。这导致了一些「显然」的一致性条件，确保引入后立即消除不会改变结果：
\begin{gather*}
  L = L \circ I_{\cat{D}} \to L \circ R \circ L \to I_{\cat{C}} \circ L = L \\
  R = I_{\cat{D}} \circ R \to R \circ L \circ R \to R \circ I_{\cat{C}} = R
\end{gather*}
这些被称为\newterm{三角恒等式}，因为它们使以下图表交换：

\begin{figure}[H]
  \centering

  \begin{subfigure}
    \centering
    \begin{tikzcd}[column sep=large, row sep=large]
      L \arrow[rd, equal] \arrow[r, "L \circ \eta"]
      & L \circ R \circ L \arrow[d, "\varepsilon \circ L"] \\
      & L
    \end{tikzcd}
  \end{subfigure}%
  \hspace{1cm}
  \begin{subfigure}
    \centering
    \begin{tikzcd}[column sep=large, row sep=large]
      R \arrow[rd, equal] \arrow[r, "\eta \circ R"]
      & R \circ L \circ R \arrow[d, "R \circ \varepsilon"] \\
      & R
    \end{tikzcd}
  \end{subfigure}
\end{figure}

\noindent
这些是函子范畴中的图表：箭头是自然变换，其复合是自然变换的水平复合。在分量层面，这些恒等式变为：
\begin{gather*}
  \varepsilon_{L d} \circ L \eta_d = \id_{L d} \\
  R \varepsilon_{c} \circ \eta_{R c} = \id_{R c}
\end{gather*}
在Haskell中，单位和余单位常以不同名称出现。单位对应到Haskell中的\code{return}函数（或\code{Applicative}定义中的\code{pure}）：

\src{snippet02}
余单位则对应\code{extract}函数：

\src{snippet03}
此处，\code{m}对应$R \circ L$的（自）函子，\code{w}对应$L \circ R$的（自）函子。正如后续所见，它们分别是单子（Monad）和余单子（Comonad）定义的一部分。

如果将自函子视为容器，单位（或\code{return}）是一个多态函数，用于围绕任意类型的值创建默认容器。余单位（或\code{extract}）则执行相反操作：从容器中提取或生成单个值。

后续我们将看到：每对伴随函子定义了一个单子和余单子。反之每个单子或余单子也可分解为伴随函子对——不过这种分解并非唯一。

在Haskell中，我们大量使用单子，但很少将其分解为伴随函子对，主要因为这些函子通常会带我们离开$\Hask$范畴。

然而可以在Haskell中定义自函子的伴随关系。\code{Data.Functor.Adjunction}模块提供了部分内容：

\src{snippet04}
这个定义需要一些解释。首先，它描述了一个多参数类型类——两个参数分别为\code{f}和\code{u}。它建立了这两个类型构造器之间的\code{Adjunction}关系。

竖线后的附加条件指定了函数依赖。例如，\code{f -> u}表示\code{u}由\code{f}确定（此处类型构造器间的关系是函数）。反之，\code{u -> f}表示若已知\code{u}，则\code{f}被唯一确定。

稍后将解释为何在Haskell中可以限定右伴随\code{u}为\newterm{可表函子（representable functor）}。

\section{伴随与Hom集}

伴随还可以通过Hom集的自然同构给出等价定义。这种定义方式与我们之前研究的泛构造紧密关联。每当你听到"存在某个唯一的态射分解某种构造"这类陈述时，应该将其理解为从某个集合到Hom集的映射。这就是"选取唯一态射"的含义。

进一步地，分解往往可以通过自然变换来描述。分解涉及交换图——某个态射等于两个态射（因子）的复合。自然变换将态射映射为交换图。因此在泛构造中，我们从态射出发经过交换图，最终得到唯一态射。这相当于建立了态射到态射的映射，或更一般地说，不同范畴中Hom集之间的映射。如果这个映射可逆且能自然扩展到所有Hom集，则构成了伴随关系。

泛构造与伴随的主要区别在于后者是全局定义的——适用于所有Hom集。例如，通过泛构造可以定义某两个特定对象的积（即使该范畴中其他对象对不存在积）。正如我们将看到的，如果范畴中任意对象对都存在积，则可以通过伴随关系重新定义这个积。

\begin{figure}[H]
  \centering
  \includegraphics[width=0.5\textwidth]{images/adj-homsets.jpg}
\end{figure}

\noindent
以下是使用Hom集的伴随替代定义。如同之前，设有两个函子 $L \Colon \cat{D} \to \cat{C}$ 和 $R \Colon \cat{C} \to \cat{D}$。任取两个对象：源对象 $d$ 在 $\cat{D}$ 中，目标对象 $c$ 在 $\cat{C}$ 中。通过 $L$ 可以将源对象 $d$ 映射到 $\cat{C}$。现在 $\cat{C}$ 中有两个对象 $L d$ 和 $c$，它们定义了一个Hom集：
\[\cat{C}(L d, c)\]
类似地，通过 $R$ 可以将目标对象 $c$ 映射到 $\cat{D}$。现在 $\cat{D}$ 中有两个对象 $d$ 和 $R c$，它们也定义了一个Hom集：
\[\cat{D}(d, R c)\]
当存在以下自然同构时，称 $L$ 是 $R$ 的左伴随：
\[\cat{C}(L d, c) \cong \cat{D}(d, R c)\]
该同构在 $d$ 和 $c$ 中都自然成立。
自然性的含义是源对象 $d$ 可以在 $\cat{D}$ 中连续变化，目标对象 $c$ 则可以在 $\cat{C}$ 中连续变化。更精确地说，存在以下两个从 $\cat{C}$ 到 $\Set$ 的共变函子之间的自然变换 $\varphi$：
\begin{gather*}
  c \to \cat{C}(L d, c) \\
  c \to \cat{D}(d, R c)
\end{gather*}
另一个自然变换 $\psi$ 作用于以下反变函子之间：
\begin{gather*}
  d \to \cat{C}(L d, c) \\
  d \to \cat{D}(d, R c)
\end{gather*}
这两个自然变换都必须可逆。

容易证明伴随的两种定义是等价的。例如，我们可以从Hom集同构推导出单位变换：
\[\cat{C}(L d, c) \cong \cat{D}(d, R c)\]
由于该同构对任意对象 $c$ 成立，特别地对 $c = L d$ 也成立：
\[\cat{C}(L d, L d) \cong \cat{D}(d, (R \circ L) d)\]
已知左侧至少包含一个态射（恒等态射）。自然变换会将这个态射映射到 $\cat{D}(d, (R \circ L) d)$ 中的元素，即插入恒等函子 $I$ 后：
\[\cat{D}(I d, (R \circ L) d)\]
我们得到了由 $d$ 参数化的态射族。它们构成了函子 $I$ 与 $R \circ L$ 之间的自然变换（自然条件容易验证），这就是我们的单位 $\eta$。

反过来，从单位和余单位的存在性也可以定义Hom集之间的变换。例如，任取Hom集 $\cat{C}(L d, c)$ 中的态射 $f$，我们要定义一个 $\varphi$ 将其映射到 $\cat{D}(d, R c)$ 中的态射。

这里的选择空间不大。一种方法是用 $R$ 提升 $f$，得到从 $R (L d)$ 到 $R c$ 的态射 $R f$，属于 $\cat{D}((R \circ L) d, R c)$。

为了获得 $\varphi$ 的分量，我们需要从 $d$ 到 $R c$ 的态射。这不是问题，因为我们可以利用 $\eta_d$ 的分量将 $d$ 映射到 $(R \circ L) d$，从而得到：
\[\varphi_f = R f \circ \eta_d\]
反向推导是类似的，$\psi$ 的推导也是如此。

回到Haskell中的 \code{Adjunction} 定义，自然变换 $\varphi$ 和 $\psi$ 分别被多态函数 \code{leftAdjunct} 和 \code{rightAdjunct}（关于 \code{a} 和 \code{b} 多态）所替代。函子 $L$ 和 $R$ 分别称为 \code{f} 和 \code{u}：

\src{snippet05}
\code{unit}/\code{counit} 表述与 \code{leftAdjunct}/\allowbreak\code{rightAdjunct} 表述之间的等价性通过这些映射体现：

\src{snippet06}
跟随范畴论描述到Haskell代码的转换过程非常具有启发性。我强烈建议作为练习进行实践。

现在我们可以解释为什么在Haskell中右伴随自动是一个\hyperref[representable-functors]{可表函子}。原因在于，粗略来说我们可以将Haskell类型范畴视为集合范畴。

当右范畴 $\cat{D}$ 是 $\Set$ 时，右伴随 $R$ 是从 $\cat{C}$ 到 $\Set$ 的函子。这样的函子是可表的，如果我们能找到 $\cat{C}$ 中的对象 $\mathit{rep}$ 使得Hom函子 $\cat{C}(\mathit{rep}, \_)$ 自然同构于 $R$。事实上，如果 $R$ 是某个从 $\Set$ 到 $\cat{C}$ 的函子 $L$ 的右伴随，这样的对象总是存在——它就是单元素集合 $()$ 在 $L$ 下的像：
\[\mathit{rep} = L ()\]
确实，伴随关系告诉我们以下两个Hom集自然同构：
\[\cat{C}(L (), c) \cong \Set((), R c)\]
对于给定的 $c$，右侧是从单元素集合 $()$ 到 $R c$ 的函数集。我们之前已经看到每个这样的函数都会从 $R c$ 中选取一个元素。这些函数的集合与 $R c$ 集合同构。因此我们有：
\[\cat{C}(L (), -) \cong R\]
这表明 $R$ 确实是可表的。

\section{从伴随来的乘积}

我们之前已经通过泛构造介绍了几个概念。当这些概念全局定义时，使用伴随更容易表达。最简单而非平凡的例子就是乘积。乘积的\hyperref[products-and-coproducts]{泛构造}本质在于：任何类乘积候选对象都可以通过泛乘积进行因子分解。

更精确地说，两个对象$a$和$b$的乘积是对象$(a\times{}b)$（在Haskell中记作\code{(a, b)}），配备两个态射$\mathit{fst}$和$\mathit{snd}$，使得对于任何其他候选对象$c$及其配备的两个态射$p \Colon c \to a$和$q \Colon c \to b$，存在唯一的态射$m \Colon c \to (a, b)$，使得$p$和$q$通过$\mathit{fst}$和$\mathit{snd}$实现因子分解。

正如我们先前所见，在Haskell中可以实现一个\code{factorizer}函数，它根据两个投影生成该态射：

\src{snippet07}
验证因子分解条件成立非常简单：

\src{snippet08}
我们有一个映射，它将一对态射\code{p}和\code{q}生成另一个态射\code{m = factorizer p q}。

如何将其转化为定义伴随所需的两个同态集之间的映射？诀窍在于跳出$\Hask$范畴，将这对态射视为积范畴中的单个态射。

让我提醒你什么是积范畴。给定任意两个范畴$\cat{C}$和$\cat{D}$，其积范畴$\cat{C}\times{}\cat{D}$的对象是对象对，分别来自$\cat{C}$和$\cat{D}$；态射也是态射对，分别来自$\cat{C}$和$\cat{D}$。

要在某个范畴$\cat{C}$中定义乘积，我们需要从积范畴$\cat{C}\times{}\cat{C}$开始。来自$\cat{C}$的态射对正是积范畴$\cat{C}\times{}\cat{C}$中的单个态射。

\begin{figure}[H]
  \centering
  \includegraphics[width=0.5\textwidth]{images/adj-productcat.jpg}
\end{figure}

\noindent
起初用积范畴来定义乘积可能会让人困惑。但这是两种截然不同的"乘积"。定义积范畴不需要泛构造——我们只需要对象对和态射对的概念即可。

然而，来自$\cat{C}$的对象对$\langle a, b \rangle$并不是$\cat{C}$中的对象，而是积范畴$\cat{C}\times{}\cat{C}$中的对象。而通过泛构造定义的$a\times{}b$（Haskell中的\code{(a, b)}）却是$\cat{C}$中的对象，它以特定方式代表了对象对$\langle a, b \rangle$。这个对象并不总是存在，即使对某些对象对存在，也可能对其他对象对不存在。

现在让我们将\code{factorizer}视为同态集之间的映射。第一个同态集位于积范畴$\cat{C}\times{}\cat{C}$中，第二个位于$\cat{C}$中。积范畴中的普通态射是一对态射$\langle f, g \rangle$：
\begin{gather*}
  f \Colon c' \to a \\
  g \Colon c'' \to b
\end{gather*}
其中$c''$可能不同于$c'$。但为了定义乘积，我们关注的是积范畴中特殊的态射——共享相同源对象$c$的$p$和$q$这对态射。这没问题：在伴随定义中，左侧同态集的源对象不是任意对象，而是左函子$L$作用于右侧范畴某对象的结果。这里合适的函子很容易猜到——就是从$\cat{C}$到$\cat{C}\times{}\cat{C}$的对角函子$\Delta$，其在对象上的作用为：
$$\Delta c = \langle c, c \rangle$$
因此伴随中的左侧同态集应该是：
$$(\cat{C}\times{}\cat{C})(\Delta c, \langle a, b \rangle)$$
这是积范畴中的同态集，其元素是我们识别为\code{factorizer}参数的态射对：
$$(c \to a) \to (c \to b) \ldots$$
右侧同态集位于$\cat{C}$中，它连接源对象$c$和某个函子$R$作用于$\cat{C}\times{}\cat{C}$中目标对象的结果。这个函子将对象对$\langle a, b \rangle$映射到我们的乘积对象$a\times{}b$。我们将这个同态集的元素视为\code{factorizer}的\emph{结果}：
$$\ldots{} \to (c \to (a, b))$$

\begin{figure}[H]
  \centering
  \includegraphics[width=0.5\textwidth]{images/adj-product.jpg}
\end{figure}

\noindent
但我们尚未建立完整的伴随关系。为此，首先需要\code{factorizer}可逆——我们在构建同态集之间的\emph{同构}。逆向的\code{factorizer}应该从某个态射$m$（即从对象$c$到乘积对象$a\times{}b$的态射）出发。换句话说，$m$应属于：
$$\cat{C}(c, a\times{}b)$$
逆向因子化器应将$m$映射到积范畴$\cat{C}\times{}\cat{C}$中的态射对$\langle p, q \rangle$，该态射对从$\langle c, c \rangle$到$\langle a, b \rangle$，即属于：
$$(\cat{C}\times{}\cat{C})(\Delta\ c, \langle a, b \rangle)$$
如果这种映射存在，我们就可以得出结论：对角函子存在右伴随。这个函子定义了乘积。

在Haskell中，我们总能通过将\code{m}分别与\code{fst}和\code{snd}组合来构造逆向的\code{factorizer}。

\begin{snip}{haskell}
p = fst . m
q = snd . m
\end{snip}
要完成这两种乘积定义方式等价性的证明，还需要证明同态集之间的映射在$a$、$b$和$c$上都是自然的。我将把这个作为练习留给认真的读者。

总结我们所做的：范畴论中的乘积可以全局地定义为对角函子的\newterm{右伴随}：
$$(\cat{C}\times{}\cat{C})(\Delta c, \langle a, b \rangle) \cong \cat{C}(c, a\times{}b)$$
这里$a\times{}b$是右伴随函子$\mathit{Product}$作用于对象对$\langle a, b \rangle$的结果。注意任何从$\cat{C}\times{}\cat{C}$出发的函子都是双函子，因此$\mathit{Product}$是双函子。在Haskell中，$\mathit{Product}$双函子简单写作\code{(,)}。你可以将它应用于两个类型来得到它们的乘积类型，例如：

\src{snippet09}

\section{从伴随来的指数对象}

指数对象 $b^a$ 或函数对象 $a \Rightarrow b$ 可以通过一个\hyperref[function-types]{普遍构造}来定义。若该构造对所有对象对都存在，则可视为一种伴随关系。此时的关键在于关注如下陈述：

\begin{quote}
对于任意其他对象 $z$ 及态射 $g \Colon z\times{}a \to b$，存在唯一的态射 $h \Colon z \to (a \Rightarrow b)$
\end{quote}
此陈述建立了同态集合之间的映射。

在此情形下，我们处理的是同一范畴中的对象，因此这对伴随函子都是自函子。左伴随函子 $L$ 作用于对象 $z$ 时产生积 $z\times{}a$，即对应与固定对象 $a$ 取积的函子。

右伴随函子 $R$ 作用于对象 $b$ 时产生函数对象 $a \Rightarrow b$（或记作 $b^a$），其中 $a$ 是固定的。这对函子间的伴随关系常写作：
\[-\times{}a \dashv (-)^a\]
通过重绘普遍构造中使用的图示，可以更直观地理解该伴随关系所蕴含的同态集合映射：

\begin{figure}[H]
  \centering
  \includegraphics[width=0.4\textwidth]{images/adj-expo.jpg}
\end{figure}

\noindent
注意 $\mathit{eval}$ 态射\footnote{参见第9章关于\hyperref[function-types]{普遍构造}} 实质上就是该伴随的余单位：
\[(a \Rightarrow b)\times{}a \to b\]
其中：
\[(a \Rightarrow b)\times{}a = (L \circ R) b\]
此前已提及，普遍构造定义的对象在同构意义下唯一。这正是为何我们称"the"积和"the"指数对象。这一性质同样适用于伴随：若某函子具有伴随，则其伴随在同构意义下唯一。

\section{挑战}

\begin{enumerate}
  \tightlist
  \item
        推导出$\psi$的自然性方阵，该变换处于以下两个(逆变)函子之间：
        \begin{gather*}
          a \to \cat{C}(L a, b) \\
          a \to \cat{D}(a, R b)
        \end{gather*}
  \item
        从伴随第二定义中的同调集同构出发推导余单位$\varepsilon$。
  \item
        完成伴随两种定义等价性的证明。
  \item
        通过余积的分解器定义出发，展示余积可通过伴随来定义。
  \item
        展示余积是对角函子的左伴随。
  \item
        在Haskell中定义乘积与函数对象之间的伴随关系。
\end{enumerate}